\documentclass[journal,onecolumn]{IEEEtran}
\usepackage{amsmath}
% Supplementary Material. Numbers moved from the main text after the third external review; sources as
% in main.tex: directive 018 (results/fm/frontier.json at 4a4ee74), directive 019 (results/hd at
% 5edad1d), directive 021 (results/hc at cac17c2).
\begin{document}
\title{Supplementary Material: Threshold Phases and the Limits of Motion Compensation for Event Streams}
\maketitle

\section*{S1. Clipping of Counts}
Counts are clipped at $15$ before coding. On one-pixel voxels with $k\ge3$, clipping removes less than
$1\%$ of the events. At $k=1$ it removes up to $6\%$ and on two-pixel voxels up to $28\%$, both by
night. On the grid of two pixels and $T_0/3$ the clip costs up to $0.16$ of fidelity at night and at
most $0.04$ of box fidelity.

\section*{S2. Scaled Luma}
In the lossy modes the luma is $16$ times the count. Coding this scaled luma without loss costs $1.9$
to $2.5$ times the bits of the unscaled lossless mode.

\section*{S3. Operating Points on a Coarse Grid}
On the grid of two pixels and $T_0/3$, x265 by day reaches $F=0.975$ with $3.26$ bits per event at
QP~$14$, $0.933$ with $1.09$ at QP~$28$, $0.840$ with $0.31$ at QP~$36$ and $0.706$ with $0.051$ at
QP~$46$. On this grid the intra picture has the default QP offset of x265.
\end{document}
